\section{Weil divisors}
\label{section-Weil-divisors}

\noindent
We will introduce Weil divisors and rational equivalence of Weil
divisors for locally Noetherian integral schemes.
Since we are not assuming our schemes are quasi-compact we have
to be a little careful when defining Weil divisors. We have to allow
infinite sums of prime divisors because a rational function may have
infinitely many poles for example. For quasi-compact schemes our
Weil divisors are finite sums as usual. Here is a basic lemma we will
often use to prove collections of closed subschemes are locally finite.

\begin{lemma}
\label{lemma-components-locally-finite}
Let $X$ be a locally Noetherian scheme. Let $Z \subset X$ be a closed
subscheme. The collection of irreducible components of $Z$
is locally finite in $X$.
\end{lemma}

\begin{proof}
Let $U \subset X$ be a quasi-compact open subscheme. Then $U$ is a Noetherian
scheme, and hence has a Noetherian underlying topological space
(Properties, Lemma \ref{properties-lemma-Noetherian-topology}).
Hence every subspace is Noetherian and
has finitely many irreducible components
(see Topology, Lemma \ref{topology-lemma-Noetherian}).
\end{proof}

\noindent
Recall that if $Z$ is an irreducible closed subset of a scheme $X$,
then the codimension of $Z$ in $X$ is equal to the dimension
of the local ring $\mathcal{O}_{X, \xi}$, where $\xi \in Z$
is the generic point. See
Properties, Lemma \ref{properties-lemma-codimension-local-ring}.

\begin{definition}
\label{definition-Weil-divisor}
Let $X$ be a locally Noetherian integral scheme.
\begin{enumerate}
\item A {\it prime divisor} is an integral closed subscheme $Z \subset X$
of codimension $1$.
\item A {\it Weil divisor} is a formal sum $D = \sum n_Z Z$ where
the sum is over prime divisors of $X$ and the collection
$\{Z \mid n_Z \not = 0\}$ is locally finite
(Topology, Definition \ref{topology-definition-locally-finite}).
\end{enumerate}
The group of all Weil divisors on $X$ is denoted $\text{Div}(X)$.
\end{definition}

\noindent
Our next task is to define the Weil divisor associated to a rational
function. In order to do this we use the order of vanishing of a
rational function along a prime divisor which is defined as follows.

\begin{definition}
\label{definition-order-vanishing}
Let $X$ be a locally Noetherian integral scheme. Let $f \in R(X)^*$.
For every prime divisor $Z \subset X$ we define the
{\it order of vanishing of $f$ along $Z$} as the integer
$$
\text{ord}_Z(f) = \text{ord}_{\mathcal{O}_{X, \xi}}(f)
$$
where the right hand side is the notion of
Algebra, Definition \ref{algebra-definition-ord}
and $\xi$ is the generic point of $Z$.
\end{definition}

\noindent
Note that for $f, g \in R(X)^*$ we have
$$
\text{ord}_Z(fg) = \text{ord}_Z(f) + \text{ord}_Z(g).
$$
Of course it can happen that $\text{ord}_Z(f) < 0$.
In this case we say that $f$ has a {\it pole} along $Z$
and that $-\text{ord}_Z(f) > 0$ is the
{\it order of pole of $f$ along $Z$}. It is important to note
that the condition $\text{ord}_Z(f) \geq 0$ is {\bf not} equivalent
to the condition $f \in \mathcal{O}_{X, \xi}$ unless the local
ring $\mathcal{O}_{X, \xi}$ is a discrete valuation ring.

\begin{lemma}
\label{lemma-divisor-locally-finite}
Let $X$ be a locally Noetherian integral scheme. Let $f \in R(X)^*$.
Then the collections
$$
\{Z \subset X \mid Z\text{ a prime divisor with generic point }\xi
\text{ and }f\text{ not in }\mathcal{O}_{X, \xi}\}
$$
and
$$
\{Z \subset X \mid Z \text{ a prime divisor and }\text{ord}_Z(f) \not = 0\}
$$
are locally finite in $X$.
\end{lemma}

\begin{proof}
There exists a nonempty open subscheme $U \subset X$ such that $f$
corresponds to a section of $\Gamma(U, \mathcal{O}_X^*)$. Hence the
prime divisors which can occur in the sets of the lemma are all
irreducible components of $X \setminus U$.
Hence Lemma \ref{lemma-components-locally-finite} gives the desired result.
\end{proof}

\noindent
This lemma allows us to make the following definition.

\begin{definition}
\label{definition-principal-divisor}
Let $X$ be a locally Noetherian integral scheme. Let $f \in R(X)^*$.
The {\it principal Weil divisor associated to $f$} is the Weil divisor
$$
\text{div}(f) = \text{div}_X(f) = \sum \text{ord}_Z(f) [Z]
$$
where the sum is over prime divisors and $\text{ord}_Z(f)$ is as in
Definition \ref{definition-order-vanishing}. This makes sense
by Lemma \ref{lemma-divisor-locally-finite}.
\end{definition}

\begin{lemma}
\label{lemma-div-additive}
Let $X$ be a locally Noetherian integral scheme. Let $f, g \in R(X)^*$. Then
$$
\text{div}_X(fg) = \text{div}_X(f) + \text{div}_X(g)
$$
as Weil divisors on $X$.
\end{lemma}

\begin{proof}
This is clear from the additivity of the $\text{ord}$ functions.
\end{proof}

\noindent
We see from the lemma above that the collection of principal Weil divisors
form a subgroup of the group of all Weil divisors. This leads to the following
definition.

\begin{definition}
\label{definition-class-group}
Let $X$ be a locally Noetherian integral scheme. The
{\it Weil divisor class group} of $X$ is the quotient of
the group of Weil divisors by the subgroup of principal Weil divisors.
Notation: $\text{Cl}(X)$.
\end{definition}

\noindent
By construction we obtain an exact complex
\begin{equation}
\label{equation-Weil-divisor-class}
R(X)^* \xrightarrow{\text{div}} \text{Div}(X) \to \text{Cl}(X) \to 0
\end{equation}
which we can think of as a presentation of $\text{Cl}(X)$. Our next task
is to relate the Weil divisor class group to the Picard group.





